% Generated by code/08_generate_results_tables.R
{\small
\setlength{\tabcolsep}{4pt}
\begin{longtable}{@{}L{1.7cm}L{7.1cm}L{4.8cm}@{}}
\caption[Hypothesis summary]{Hypothesis summary}
\label{tab:hypothesis_result_summary}\\
\toprule
Hypothesis & Main result & Interpretation \\
\midrule
\endfirsthead
\multicolumn{3}{c}{\textit{Table~\ref{tab:hypothesis_result_summary}---continued}} \\
\toprule
Hypothesis & Main result & Interpretation \\
\midrule
\endhead
\midrule
\multicolumn{3}{r}{\textit{Continued on next page}}\\
\endfoot
\bottomrule
\endlastfoot
H1 & The distributional branch raises the EU27 Stage 1 liability rate relative to the jurisdiction-average rule in every year from 2017 to 2022. & Supported. \\
H2 & ETR measurement explains 69.1\% of pooled scenario variance. & Supported. \\
H3 & The SbS regime explains 15.0\% of pooled scenario variance and reduces preferred-core EU revenue by US\$138.6 bn through the UTPR channel. & Supported. SbS is the second-largest variance source and materially changes collection channels. \\
H4 & Tax terms are jointly significant; baseline BIC selects Linear; logarithmic 5\% semi-elasticity is 2.00. & Supported in its cautious form: the panel supports a negative tax-profit relationship, but not a sharply nonlinear lower-tail specification. \\
H5 & PC1 explains 48.8\% of typology variation and the selected four-cluster solution has 100.0\% jitter stability. & Supported. Member States separate into high-salience conduit hubs, QDMTT-protected smaller exposures, residual IIR/UTPR exposure states and zero-current-exposure states. \\
\end{longtable}
}

\vspace{0.35em}
\noindent\begin{minipage}{\textwidth}
\footnotesize\raggedright\textit{Notes:} The table summarises the evidence reported in the Stage 1 to Stage 4.6 diagnostics and the country fiscal-exposure calculations in Section \ref{sec:results_sbs}.
\end{minipage}
